\section{从功能渗透到Tier-1收入：链条的正确读法}

智能汽车电子的长期方向是电子电气架构从大量分布式ECU向座舱域控、驾驶域控、中央计算和更强的软件协同迁移。对OEM而言，功能组合影响车型竞争力；对Tier-1而言，真正的财务链条仍是定点、开发验证、SOP、批量交付、客户验收、收入确认、回款。行业景气可以推动前端需求，却不能压缩中间的开发、认证、价格和现金环节。

德赛西威的三条产品线把这个路径具体化。智能座舱包含信息娱乐、显示、HUD和座舱域控；智能驾驶包含域控、传感器和算法；网联服务则覆盖T-Box、OTA、基础软件和智能进入。系统化产品可以提高单车配置价值，但同时带来高算力SoC、显示/感知BOM、功能安全和研发投入的约束。

\begin{exhibitbox}[2025年已确认的产品池]
\small
\begin{tabularx}{\textwidth}{L{2.6cm}Y Y Y}
\toprule
产品池 & 收入与增速 & 毛利率与质量 & 估值含义 \\
\midrule
智能座舱 & CNY20.585bn，占63.23\%，同比+12.92\% & 18.83\%，为最大已确认收入/毛利池 & 提供稳定盈利基础；客户份额和ASP未披露 \\
智能驾驶 & CNY9.700bn，占29.79\%，同比+32.63\% & 16.36\%，同比-3.55pct & 代表增长，但不能假定综合毛利扩张 \\
网联及其他 & CNY2.272bn，占6.98\%，同比+9.52\% & 分部毛利未披露 & 有服务基础，但不给纯软件估值溢价 \\
\bottomrule
\end{tabularx}
\sourcenote{公司2025年年报。}
\end{exhibitbox}

上表的投资含义并不是“智驾收入快，所以应当给更高的倍数”。2025年智能驾驶成本增长快于收入，毛利率低于座舱；若SoC和传感器成本、研发投入或OEM议价没有改善，组合向智驾迁移可带来收入增长，却未必带来同比例的净利增长。基准估值因此用合并盈利和P/E，而不是把智驾板块按软件或纯算法公司估值。

\section{技术架构与价值量的边界}

公司披露第五代AI座舱、端云一体方案、中央计算、3D/4D毫米波雷达和舱驾融合等产品进展。技术能力可以解释为什么公司仍具备项目竞争力；但中央计算只有“国内多家头部客户订单”的匿名披露，AR-HUD、区域控制器有订单或即将量产描述，均缺少客户金额、车型、SOP、ASP、产品毛利和收入确认。它们应当被记为可选性，而不是基准利润。

\begin{keyinsight}[技术论点的财务门槛]
对于每一项新产品，估值信用须通过四道门：客户/平台确认、SOP或量产、价格或价值量、利润/现金转换。只有第一道门的产品进展可以说明技术方向；若缺后二至四道门，报告只能把它写为研究跟踪项。
\end{keyinsight}

座舱业务的第四代平台已经在理想、小米、吉利、奇瑞和广汽等客户中配套，提供较高质量的量产事实。智能驾驶全系列平台也有多OEM大规模量产的年报描述。二者构成公司“系统交付”定位的实证，但不允许把客户名单按车型销量或品牌规模分配成公司收入。

\section{行业需求与公司利润池的错位}

行业需求可以被销量、NOA渗透率、多屏/算力配置和OEM平台换代验证；公司利润池还要看BOM、价格竞争、研发、工厂利用率和回款。2025年公司销售量40.346m套、生产量同比增长39.71\%，支持规模扩张；年末库存和应收规模也提醒读者，交付增长的质量必须由现金流和减值风险共同确认。

因此本报告把未来行业空间作为方向性背景，而把已确认分部收入、分部毛利、原材料占成本、应收/库存/合同负债和经营现金流作为基础估值的输入。这个顺序减少了主题热度被误写为EPS的风险。
